% !TeX root = ../../Book2.tex
\section{Fourier 变换法}
\label{b2:sec:1202}

在全空间中，空间平移不会改变方程，离散特征级数因而换成连续频率积分。变换负责找出核，卷积与逼近恒等负责证明粗初值确实演化成解。外力也可逐时叠加，但积分公式本身并不自动赋予它任意阶的经典导数。

\subsection{热核}
\label{b2:subsec:120201}

全空间没有端点，也没有正的最低允许频率。Fourier 变换把每个连续频率单独演化。先对足够光滑且衰减的解形式计算，按照本卷的 \(2\pi\) 约定有
\[
 \partial_t\widehat u(t,\xi)
 =-4\pi^2\kappa|\xi|^2\widehat u(t,\xi),
 \quad\widehat u(t,\xi)=e^{-4\pi^2\kappa t|\xi|^2}\widehat f(\xi).
\]
\cref{b2:thm:03-gaussian}已经计算了这个乘子的逆变换。沿用第二编的热核记号，在 \(\mathbb R^d\) 上写作
\begin{equation}\label{b2:eq:12-gauss-heat}
 H_t(x)=(4\pi t)^{-d/2}e^{-|x|^2/(4t)},
 \quad t>0,
 \quad\widehat H_t(\xi)=e^{-4\pi^2t|\xi|^2}.
\end{equation}
方程扩散系数为 \(\kappa\) 时使用 \(H_{\kappa t}\)，不另换 Fourier 归一化。这个核非负、总积分为一；其空间宽度约为 \(\sqrt t\)，高度约为 \(t^{-d/2}\)。由直接微分有 \(\partial_tH_{\kappa t}=\kappa\Delta H_{\kappa t}\)。

\ProseParagraphBreak

这些计算可以独立验证卷积候选，不要求先假定未知解的变换存在。对于粗初值，求解的逻辑次序是先用核定义 \(u\)，再证明可微性和初值收敛；前面的频率微分只负责找到候选。第二编已经证明 Gauss 核的逼近恒等性质，下面把其参数取为 \(\sqrt{4\pi\kappa t}\)。

\subsection{卷积表示}
\label{b2:subsec:120202}

\Needspace{6\baselineskip}
\begin{theorem}[全空间热核解]\label{b2:thm:12-heat-convolution}
设整数 \(d\geq1\)、\(\kappa>0\)，\(H_s(x)=(4\pi s)^{-d/2}e^{-|x|^2/(4s)}\)（\(s>0\)、\(x\in\mathbb R^d\)）。设 \(1\leq p\leq\infty\)、\(f\in L^p(\mathbb R^d;\mathbb C)\)。则
\begin{equation}\label{b2:eq:12-heat-convolution}
 u(t,x)=\int_{\mathbb R^d}H_{\kappa t}(x-y)\cdot f(y)\,\mathrm{d}y
\end{equation}
对 \(t>0\)、\(x\in\mathbb R^d\) 有意义，且在 \((0,\infty)\times\mathbb R^d\) 上光滑，满足 \(u_t=\kappa\Delta u\)。当 \(t\downarrow0\) 且 \(p<\infty\) 时，\(u(t)\to f\) 于 \(L^p\)。若选定的代表 \(\widetilde f\) 有界一致连续，则 \(u(t)\to\widetilde f\) 在一致范数中成立；若选定 \(f\) 的一个有界代表 \(\widetilde f\)，且它在 \(x_0\in\mathbb R^d\) 连续，则 \(u(t,x_0)\to\widetilde f(x_0)\)。
\end{theorem}
\begin{proof}
热核及其每阶时空导数属于每个 \(L^r\)、\(1\leq r\leq\infty\)。固定 \(t\) 的正紧区间与 \(x\) 的紧集后，每个导数可由一个平移后仍快速衰减的函数统一控制。对 \(f\in L^p\)，Hölder 不等式给出绝对可积性；对差商使用同一控制，可以把每一阶导数移入积分。因此卷积光滑，且核的微分等式传给 \(u\)。这里在 \(p=1\) 用核导数的一致界，在 \(p=\infty\) 用其绝对可积性，两个端点都被覆盖。

有限 \(p\) 与一致连续代表的情形由\cref{b2:thm:03-approximate-identity}得到。卷积不因初值在零测集上的修改而改变。对选定代表的连续点，把 \(\widetilde f(x_0-y)-\widetilde f(x_0)\) 的积分分成 \(|y|<\delta\) 与补集。近处由连续性控制，远处由 \(2\sup_x|\widetilde f(x)|\) 乘热核尾质量控制，先选 \(\delta\) 再令 \(t\downarrow0\) 即得。这样取得的点值是所选连续代表的点值，而不是给 \(L^p\) 等价类额外指定一个任意值。
\end{proof}

初值也可以是一点质量，此时解是 \(H_{\kappa t}(x-y_0)\)，在 \(t\downarrow0\) 时只以测度或分布意义趋于 \(\delta_{y_0}\)。为了区别空间核与时空基本解，令 \(E(t,x)=H_{\kappa t}(x)\) 当 \(t>0\)，在 \(t<0\) 令 \(E=0\)。那么
\[
 (\partial_t-\kappa\Delta_x)E=\delta_{(0,0)}
\]
在时空分布中成立。确实，对紧支集光滑 \(\varphi(t,x)\)，从 \(t=\varepsilon\) 起分部积分，热方程抵消内部项，留下 \(\int H_{\kappa\varepsilon}(x)\cdot\varphi(\varepsilon,x)\,\mathrm{d}x\)。它由逼近恒等趋于 \(\varphi(0,0)\)。同时 \(E\) 在有限时间带的空间积分为该带长度，所以零时附近的局部可积性也已得到。

\ProseParagraphBreak

对非齐次方程，设整数 \(d\geq1\)、\(\kappa,T>0\)、\(1\leq p<\infty\)，\(f\in L^p(\mathbb R^d;\mathbb C)\)、\(F\in C([0,T];L^p(\mathbb R^d;\mathbb C))\)。由逐时叠加得到
\begin{equation}\label{b2:eq:12-duhamel}
 u(t)=H_{\kappa t}*f+
 \int_0^tH_{\kappa(t-s)}*F(s)\,\mathrm{d}s.
\end{equation}
在积分的右端按 \(H_0*F(t)=F(t)\) 理解。被积函数是连续的 \(L^p\) 值函数，积分可作为 Riemann 和的范数极限；Young 不等式给出
\[
 \left\|\int_0^tH_{\kappa(t-s)}*F(s)\,\mathrm{d}s\right\|_p
 \leq\int_0^t\|F(s)\|_p\,\mathrm{d}s.
\]
这保证源项趋于零初值。先对光滑紧支集的时空源项求导，再以连续 \(L^p\) 范数逼近一般 \(F\)，便在分布中得到 \(u_t-\kappa\Delta u=F\)。逼近可将时间分成有限小段，对有限个空间取值分别磨光，再以时间连续插值连接；上式使解同时收敛，测试函数的积分配对也连续。对于这里的一般连续 \(L^p\) 值源项，方程按分布意义成立；经典二阶空间正则性还需进一步的条件。

\subsection{高维热方程}
\label{b2:subsec:120203}

高维热核是各坐标一维核的乘积，但高维扩散的幅度不能从一维直观中省略维数。换元 \(x=\sqrt t\,z\) 给出
\[
 \|H_t\|_r=C_{d,r}t^{-\frac d2(1-1/r)},
 \quad1\leq r\leq\infty.
\]
例如有限质量初温满足 \(\|u(t)\|_\infty\leq(4\pi\kappa t)^{-d/2}\|f\|_1\)。没有任何质量损失也能使温度峰值降低，因为同一热量散布到越来越大的区域。

\ProseParagraphBreak

方向上的扩散速率也可以不同。若整数 \(d\geq1\)，\(B\) 是实对称正定的 \(d\times d\) 常矩阵，方程为 \(u_t=\sum_{i,j=1}^dB_{ij}\partial_i\partial_ju\)，则 \(t>0\) 时的相应核是
\[
 H_t^B(x)=\dfrac1{(4\pi t)^{d/2}\sqrt{\det B}}
 \exp\left(-\dfrac{x^{\mathsf T}B^{-1}x}{4t}\right).
\]
正交对角化 \(B\) 后，各方向分别使用一维扩散核，换元的行列式给出归一化因子。故核仍非负、质量为一，变换为 \(e^{-4\pi^2t\xi^{\mathsf T}B\xi}\)。正定性使每个方向都发生扩散。若 \(B\) 仅半正定且有零特征值，\(B^{-1}\) 与 \((\det B)^{-1/2}\) 就没有上述意义，不能在函数密度公式中直接把某个特征值置零。

\ProseParagraphBreak

例如在 \(\mathbb R^2\) 上取 \(B=\operatorname{diag}(1,0)\)，方程成为 \(u_t=u_{xx}\)，没有 \(y\) 方向的扩散。此时传播核是乘积测度
\[
 H_t^{(1)}(x)\otimes\delta_0(y),\qquad
 H_t^{(1)}(x)=(4\pi t)^{-1/2}e^{-x^2/(4t)},\quad t>0,
\]
它支撑在直线 \(y=0\) 上，不是相对于二维 Lebesgue 测度的可积函数密度。对 \(f\in L^2(\mathbb R^2;\mathbb C)\)，传播只沿 \(x\) 作卷积：
\[
 u(t,x,y)=\int_{\mathbb R}H_t^{(1)}(x-z)f(z,y)\,\mathrm{d}z.
\]
Fubini 定理使几乎每个 \(y\) 对应的初值切片属于 \(L^2(\mathbb R_x)\)，故这一维卷积有意义，并在分布意义下满足 \(u_t=u_{xx}\)。但它不提高横向正则性：取非负非零 \(\phi\in C_c^\infty(\mathbb R;\mathbb R)\)，令 \(f(x,y)=\phi(x)\mathbf1_{(0,1)}(y)\)，就有
\[
 u(t,x,y)=\bigl(H_t^{(1)}*\phi\bigr)(x)\mathbf1_{(0,1)}(y).
\]
前一个因子在 \(t>0\) 时光滑且处处为正，\(y=0,1\) 处的跳跃却仍保留。一般半正定矩阵经正交对角化后也如此：正特征值方向使用 Gaussian 核，零特征值方向保留点质量，不能从正方向的磨光推出全空间光滑化。

\ProseParagraphBreak

这一正定二次型公式可与 John K. Hunter、Bruno Nachtergaele 的《Applied Analysis》\cite[pp. 303--304]{Hunter2001Applied}对照，其变换采用不同的 \(2\pi\) 归一化。本卷已经在\secref{b2:sec:0502}换算过常数，这里直接使用一致的频率约定。
